\documentclass[12pt]{article}

\usepackage[utf8]{inputenc}
\usepackage[spanish]{babel}
\usepackage[table, dvipsnames]{xcolor}
\usepackage[hidelinks]{hyperref}
\decimalpoint

\usepackage[heuristica,vvarbb,bigdelims]{newtxmath}
\usepackage[T1]{fontenc}
\renewcommand*\oldstylenums[1]{\textosf{#1}}


\usepackage{csquotes} %Authors and Quotes
\usepackage{authblk}


\usepackage{array} %Tabulars and arrays
\usepackage{tabularray}
\usepackage{multirow}
\usepackage{multicol}
\usepackage{tabularx}
\usepackage{booktabs}
\usepackage{extarrows}


\usepackage[letterpaper,top=2cm,bottom=2.1cm,left=2cm,right=2cm,marginparwidth=1.75cm]{geometry}


\usepackage{physics} %Math and symbols and more of that
\usepackage{textcomp, gensymb}
\usepackage{amsmath}
\usepackage{amsbsy}
\usepackage{upgreek}
\usepackage{mathtools}
\usepackage{wasysym}


\usepackage{comment} %Styles and more
\usepackage{cancel}
\usepackage{enumerate}
\usepackage[framemethod=default]{mdframed} 
\usepackage{tikz}
\usepackage{graphicx}
\usepackage{fancyhdr}
\usepackage{mathdots}
\usepackage{color}
\usepackage{blindtext}
\usepackage{pageslts}
\usepackage{enumitem}
\usepackage{anyfontsize}
\usepackage{boldline}
\usepackage{siunitx}
\usepackage[explicit]{titlesec}


\usepackage{lipsum} %Miscellaneous
\usepackage{float}

%Bibliography data management
\usepackage{biblatex} %Imports biblatex package
% \addbibresource{citation.bib} %Import the bibliography file


\pagenumbering{arabic}

\setlength{\parindent}{0pt}

\usetikzlibrary{patterns}
\usetikzlibrary{shadows.blur}
\usetikzlibrary{shapes}

\DeclareMathOperator\arccosh{arccosh}

%Page styles and more ##########################

\mdfdefinestyle{questions}{
linecolor=BrickRed,
linewidth=2pt,
leftmargin=1cm,
rightmargin=1cm
}

\mdfdefinestyle{solutions}{
backgroundcolor = Cyan!10,
linecolor=Purple,
linewidth=2pt,
}

\mdfdefinestyle{draw}{
linecolor=Black,
linewidth=1.5pt,
}

\mdfdefinestyle{equation}{
    backgroundcolor = black!6,
    linecolor=blue,
    linewidth=1.5pt,
    topline=false,
    rightline=false,    
    leftline = false,
}

\mdfdefinestyle{membret}{
    linecolor = blue,
    linewidth = 1.5pt,
    topline = false,
    rightline = false,    
    leftline = false,
}

%KeyWords-Command ##############################
\titleformat{\section}
  {\normalfont\Large\bfseries} % Text style
  {}                           % Label (handled in the next argument)
  {0pt}                        % Separation
  {%
    \makebox[0pt][l]{\color{blue}\rule[-4pt]{4pt}{1.2em}} % The Vertical Bar
    \hspace{10pt}\thesection\quad#1                      % The Title
  }

% Format for Subsection
\titleformat{\subsection}{\normalfont\large\bfseries}
  {}
  {0pt}
  {%
    \makebox[0pt][l]{\color{black}\rule[-2pt]{3pt}{1em}} % Thinner bar for subsection
    \hspace{10pt}\thesubsection\quad#1
}


%KeyWords-Command ##############################

\def\arraystretch{1.2}
\def\QWidth{2.5pt}

\newenvironment{question}[2]
    {
    {\global\arrayrulewidth = \QWidth}
    \def\arraystretch{1.5}
    \def\points{#2} %Save the points in a variable to use later
    
    \begin{center}
    \begin{tabular}{
    !{\color{red} \vrule width \QWidth} m{\linewidth*\real{0.05}} !{\color{red} \vrule width \QWidth} m{\linewidth*\real{0.75}} !{\color{red} \vrule width \QWidth} m{\linewidth*\real{0.1}} !{\color{red} \vrule width \QWidth}}  % Begin center and table 
    
    \arrayrulecolor{red}\hline 
    
    \centering{\bf #1}
    & 
    %\vspace{2mm}\begin{minipage}{\linewidth}
    
    }{
        %\end{minipage} \vspace{-1.5mm}
        &
        {\bf +\points pts.}
        \\ \arrayrulecolor{red}\hline
        
        \end{tabular}
        \end{center}
        {\global\arrayrulewidth = 1pt}
    }


\makeatletter
\renewcommand\maketitle{
    {\centering
    \phantom{Hello there} \\
    \vspace{5cm}
    \@title \\
    \large
    \vspace{10mm}
    Miguel Doradea Meléndez \\
    Samuel Martín Arteaga \\
    \vspace{3mm}
    Universidad Politécnica de Valencia
   
    \clearpage}
}


\makeatother

\fancypagestyle{format1}{
    \renewcommand{\headrulewidth}{0pt}
    \renewcommand{\footrule}{ {\color{blue} \rule{\linewidth}{1.5pt}}}
    
    \fancyhead{}
    \fancyfoot{}
    
    \fancyhf[FC]{\hspace{7mm} GIE: EDSO \hfill Página: \thepage \ de \lastpageref{LastPages}}
}

%Format and title#######################

\pagestyle{format1}
\title{{\fontsize{32}{32}\selectfont \bf Grado de Ingeniería de la Energía} \\ \vspace{3mm}
    {\LARGE \color{Blue} \textbf{Energías de Desarrollo Sostenible:} Práctica 2, Diseño de un parque eólico y estimación de su viabilidad económica con el programa RETScreen.}}

%Begin Document####################### 
\begin{document}
\maketitle

\section{Objetivos}

El objetivo principal de la presente práctica es diseñar de manera simplificada un parque eólico. Para ello se dividirá la práctica en dos partes:

\begin{itemize}[label = \textbullet]
    \item Estudio de la metodología de diseño de un parque eólico
    \item Estudiar la viabilidad económica
\end{itemize}

\section{Metodología de Diseño de un Parque Eólico}

Para localizar una buena ubicación del parque debemos buscar una colina que sea capaz de albergar la cantidad deseada de aerogeneradores. Se desea instalar un parque de $6 \unit{MW}$, por lo que se necesitan \textbf{3 aerogeneradores de 2 MW}. Vamos a suponer que utilizamos aerogeneradores \textbf{G83-2MW} de Gamesa. En el catálogo podemos ver que la longitud de pala es de $40.5 \unit{m}$, lo que equivale a $83 \unit{m}$ de diámetro de rotor y una potencia de $2 \unit{MW}$. \medskip

La información técnica del aerogenerador escogido (\textbf{G83-2MW}) nos da la curva de Potencia del Aerogenerador:

\begin{table}[htbp]
    \centering
    \small
    % Left Minipage
    \begin{minipage}{0.45\textwidth}
        \centering
        \begin{tblr}{
            colspec = {c c},
            hlines = {0.5pt, solid, gray!50}, 
            vlines = {0.5pt, solid, gray!50},
            row{1,2} = {bg=blue!20, font=\bfseries}, 
            hline{3} = {1.5pt, solid, black},        
        }
        Velocidad del viento & Curva de potencia \\
        v(m/s)               & P(kW)             \\
        0                    & 0                 \\
        1                    & 0                 \\
        2                    & 0                 \\
        3                    & 0                 \\
        4                    & 65.1              \\
        5                    & 152               \\
        6                    & 285               \\
        7                    & 471               \\
        8                    & 716               \\
        9                    & 1025              \\
        10                   & 1377              \\
        11                   & 1691              \\
        12                   & 1882              \\
        \end{tblr}
    \end{minipage}\hfill % \hfill pushes the two minipages apart
    % Right Minipage
    \begin{minipage}{0.45\textwidth}
        \centering
        \begin{tblr}{
            colspec = {c c},
            hlines = {0.5pt, solid, gray!50}, 
            vlines = {0.5pt, solid, gray!50},
            row{1,2} = {bg=blue!20, font=\bfseries}, 
            hline{3} = {1.5pt, solid, black},        
        }
        Velocidad del viento & Curva de potencia \\
        v(m/s)               & P(kW)             \\
        13                   & 1964              \\
        14                   & 1990              \\
        15                   & 1998              \\
        16                   & 1999              \\
        17                   & 2000              \\
        18                   & 2000              \\
        19                   & 2000              \\
        20                   & 2000              \\
        21                   & 2000              \\
        22                   & 2000              \\
        23                   & 2000              \\
        24                   & 2000              \\
        25                   & 2000              \\
        \end{tblr}
    \end{minipage}
    \caption{Información técnica del aerogenerador G83-2MW de la marca Gamesa.}
    \label{tab:1}
\end{table}

La curva de potencia nos da la potencia que el aerogenerador es capaz de generar a cada velocidad del viento. Se puede observar que a partir de los $12 \unit{m/s}$ el aerogenerador alcanza su potencia nominal de $2 \unit{MW}$, por lo que no es capaz de generar más energía a velocidades del viento superiores a esta. En la siguiente gráfica podemos observar este comportamiento: \medskip



\subsection{Pasos a seguir para el planteamiento del parque eólico}

\begin{enumerate}[label = \textbf{Paso \arabic*:}, leftmargin = 2.5cm, labelindent = 2.5cm]
    \item Información necesaria para seleccionar una buena ubicación de un parque eólico.
    \item Herramientas necesarias.
    \item Elegir una zona eólica.
    \item Búsqueda de una colina.
    \item Caracterización eólica de la zona.
    \item Cálculo de la energía bruta y neta
\end{enumerate}

\subsection{Requisitos}

\subsubsection{Administrativos}

Plan Eólico de la Comunitat Valenciana.
Se debe comprobar que el parque eólico no incumpla las condiciones
mínimas estipuladas por el Plan Eólico de la Comunidad Valenciana (DOGV,
2001) y por el PER (IDAE, 2011):

\begin{itemize}[label = \textbullet]
    \item  Distancia mínima a suelo urbano o urbanizable de uso no industrial (Plan Eólico): 1000 m.
    \item  Distancia mínima a una línea de la red eléctrica de alta tensión (PER): 250 m.
    \item  Distancia mínima a carreteras autonómicas (PER): 100 m.
    \item  Distancia mínima a carreteras nacionales, autovía y/o autopista (PER): 200 m.
    \item  Para consultar el plan urbanístico del suelo se emplea la aplicación Terrasit (IDECV, 2017).
\end{itemize}

\subsubsection{Técnicos}
\begin{itemize}
    \item Orografía. Cimas largas y poco escarpadas.
    \item Dirección predominante del viento. Interesa que sea lo más perpendicular posible a la colina para que los aerogeneradores no se hagan sombra aerodinámica.
    \item Velocidad media del viento elevada.
    \item Distribución de frecuencias del viento. Weibull (C, K).
    \item Cumplir con distancias a poblaciones, carreteras, líneas transporte eléctrico (Relacionado con Requisitos Administrativos).
\end{itemize}



\section{Cálculo de la energía bruta y neta}

Para calcular la energía bruta y neta se va a utilizar una plantilla \textbf{Excel de PoliformaT (Cálculo energía.xlsx)}. \medskip

Antes de hacer uso de la Excel hay que entender la importancia del valor de la densidad para el cálculo de la Potencia del Aerogenerador

\[ P_{extraible} = \frac{1}{2} \rho \cdot A \cdot v_1^3 \cdot C_p \]

Hay que tener en cuenta el cambio de presión atmosférica a nivel del mar y a la altitud considerada:

\begin{itemize}
    \item La densidad de aire a $25\celsius$ y nivel del mar es $1.225 \unit{kg/m^3}$
    \item Entre 900 y 1100 metros de altitud (incluido el buje) y a temperaturas entre $10\celsius$ y $15 \celsius$, la densidad del aire está entorno a $1.135 \unit{kg/m^3}$
\end{itemize}

Queda por calcular el número de horas equivalentes de un parque eólico. Esto es el cociente de la producción anual de energía (PAE) y su potencia nominal.

\begin{equation}
    NHE = \frac{PAE}{P_{nom}}
\end{equation}

El \textbf{factor de carga} es el cociente entre el NHE de un año y el número de horas de un año.

\begin{equation}
    CF = \frac{NHE}{8760 \unit{h}}
\end{equation}

A partir de un factor de carga del 30\% se considera que el PE tiene posibilidades de ser rentable técnicamente. \medskip

Utilizando la hoja de cálculo de excel conseguimos la siguiente información (la hoja de excel está anexada con este documento):


\begin{table}[h!]
    \centering
    \renewcommand{\arraystretch}{1.3}
    \begin{tblr}{
        colspec = {c c c c c c c},
        row{1-2} = {c, font=\bfseries, bg=blue!20},
        hline{3} = {2pt, solid},
    }
        \textbf{Dir.} & \textbf{Sector} & \textbf{C} & \textbf{K} & \textbf{freq} & \textbf{Suma Parcial} & \textbf{Energía Bruta} \\ 
        & & \scriptsize linea-10 & \scriptsize linea-11 & \scriptsize linea-5/10000 & & \scriptsize (kWh/sector) \\
        N   & 1  & 7.44  & 1.289 & 4.06\%  & 4999025.08 & \textbf{202960.42} \\
        NNE & 2  & 2.93  & 0.909 & 2.23\%  & 1380568.17 & \textbf{30786.67} \\
        ENE & 3  & 6.38  & 1.314 & 5.36\%  & 4002826.38 & \textbf{214551.49} \\
        E   & 4  & 6.24  & 1.783 & 12.16\% & 3356322.41 & \textbf{408128.81} \\
        ESE & 5  & 4.72  & 1.786 & 10.61\% & 1598320.26 & \textbf{169581.78} \\
        SSE & 6  & 3.64  & 1.867 & 6.55\%  & 632223.27  & \textbf{41410.62} \\
        S   & 7  & 3.21  & 1.687 & 4.42\%  & 476872.44  & \textbf{21077.76} \\
        SSO & 8  & 3.15  & 0.952 & 3.53\%  & 1473506.14 & \textbf{52014.77} \\
        OSO & 9  & 7.70  & 1.169 & 7.86\%  & 5137085.32 & \textbf{403774.91} \\
        O   & 10 & 11.05 & 1.700 & 21.60\% & 8207053.37 & \textbf{1772723.53} \\
        ONO & 11 & 9.70  & 1.567 & 13.08\% & 7040036.04 & \textbf{920836.71} \\
        NNO & 12 & 10.24 & 1.680 & 9.09\%  & 7610896.34 & \textbf{691830.48} \\
    \end{tblr}
    \caption{Cálculo de la energía bruta por sector de viento.}
    \label{tab:2}
\end{table}

Cambie algunas cosas en el archivo, en teoría se debe de poder ver desde aquí

\end{document}